% Source-derived chapter 11: The induction morphism
% Original source: rigid-local-systems-multiplicative-eigenvalue-problem
\chapter{The induction morphism}
\label{rigid-local-systems-multiplicative-eigenvalue-problem-chapter-11}
\source{rigid-local-systems-multiplicative-eigenvalue-problem}

\begin{remark}[Source section]
\label{rigid-local-systems-multiplicative-eigenvalue-problem-chapter-11-source}
\source{rigid-local-systems-multiplicative-eigenvalue-problem}
\source{rigid-local-systems-multiplicative-eigenvalue-problem}
This chapter reproduces the corresponding section of the retrieved
LaTeX source, including its definitions, statements, examples, and
proofs. It has not yet been translated into Lean.
\notready
\end{remark}

% The source section title was: The induction morphism
\label{noncan}
\source{rigid-local-systems-multiplicative-eigenvalue-problem}
Our aim now is to show that  $\Pic'=\Pic'(\mf)$ and $\Pic'^{+,\deg=0}(\mf)$ (Definition  \ref{Br2} and Proposition \ref{Br4}) are both controlled by the Levi-subgroup by an induction procedure. Since we need shifting procedures in the proof, we operate in a more general setting:
Assume that the Gromov-Witten number $\langle\sigma_{I^1}, \sigma_{I^2},\dots,\sigma_{I^s}\rangle_{d,D}=1$, where $D=-\deg(\mn)$. This implies that $\Omega(d,r,\mn,n,\vec{I})\to \Par_{n,\mn,S}$ is birational, and \eqref{timer1} holds.

\begin{defi}
\source{rigid-local-systems-multiplicative-eigenvalue-problem}
\notready
Define $\Pic'=\Pic'(\mf)\subseteq\Pic(\Par_{n,\mn,S})$ as follows (generalizing  Definition  \ref{Br2} to the case of arbitrary $\mathcal{N}$): It is subgroup formed by  $\mathcal{B}(\vec{\lambda},\ell)$ such that for any $j$, both of the following two conditions hold : (a) $\lambda^j_{a-1}=\lambda^{j}_{a}$ whenever $a\in I^j$, $a>1$ and $a-1\not \in I^j$ and (b) If $n\not\in I^j$ and $1\in I^j$ then $\lambda^j_1=\lambda^{j}_{n}+\ell$.
\end{defi}
It is easy to see that $\Pic'$ behaves well under shift operations of the tuple $(d,r,\mn,n,\vec{I})$: In \eqref{basicD} with $\vec{J}=I$, the operation $\op{Sh}_p$ pulls back $\Pic'$ for the tuple $(d',r,\mn(p),n,\vec{K})$ to $\Pic'$ for $(d,r,\mn,n,\vec{I})$ (use Proposition  \ref{easily}).


Let $\Pic^{\deg=0}(\Par_{n,\mn,S})$ be the subgroup given by line bundles $\mathcal{B}(\vec{\lambda},\ell)$ such that  equality holds in the following inequality (this is the same inequality as \eqref{wallnew} when $D=0$)
\begin{equation}\label{wallneww}
\source{rigid-local-systems-multiplicative-eigenvalue-problem}
\frac{1}{r}\bigl(-d+\sum_{j=1}^s\sum_{k\in I^j} \frac{\lambda^j_k}{\ell}\bigr)\ \leq\  \frac{1}{n}\bigl(-D +\sum_{j=1}^s \frac{|\lambda^j|}{\ell}\bigr)
\end{equation}
Finally $\mf\subseteq \Pic^{+,\deg=0}_{\Bbb{Q}}(\Par_{n,\mn,S})$ is the effective $\Bbb{Q}$ subcone.
\begin{defi}
\source{rigid-local-systems-multiplicative-eigenvalue-problem}
\notready\label{Pin}
\source{rigid-local-systems-multiplicative-eigenvalue-problem}
\begin{enumerate}
\item Any line bundle $\ml$ on $\ParL$ (or on $\ParL-\mathcal{R}_L$ ) has an index $m(\mathcal{L})\in \Bbb{Z}$: $t\in\Bbb{C}^*$ acts on points $(\mv,\mq,\mf,\mg,\gamma)\in\ParL$, by multiplication
by $t^{n-r}$ on $\mv$ and $t^{-r}$ on $\mq$ and leaving $\gamma$ unchanged. This action lifts to the line bundle $\ml$, as multiplication by $t^m$, for some $m\in\Bbb{Z}$. We set $m(\mathcal{L})=m$.
\item Let $\Pic^{\deg=0}_{\Bbb{Q}}(\ParL)\subset \Pic_{\Bbb{Q}}(\ParL)$ (similarly $\Pic^{\deg=0}_{\Bbb{Q}}(\ParL-\mathcal{R}_L)\subset \Pic_{\Bbb{Q}}(\ParL-\mathcal{R}_L)$) be  the subspace of line bundles of index $0$.
\item Let $\Pic^{+,\deg=0}_{\Bbb{Q}}(\ParL)\subset \Pic_{\Bbb{Q}}(\ParL)$ (similarly $\Pic^{+deg=0}_{\Bbb{Q}}(\ParL-\mathcal{R}_L)\subset \Pic_{\Bbb{Q}}(\ParL-\mathcal{R}_L))$be  the subspace of effective line bundles (they automatically have  index $0$).
\end{enumerate}
\end{defi}
In this section, we set up an induction map
\begin{equation}\label{inducto}
\source{rigid-local-systems-multiplicative-eigenvalue-problem}
\op{Ind}:\Pic(\ParL-\mathcal{R}_L)\to \Pic(\Par_{n,\mn,S}).
\end{equation}
We then show the following properties:
\begin{enumerate}
\item [(I1)] The map $\Pic^{+,\deg=0}_{\Bbb{Q}}(\ParL)\to \Pic^{+,\deg=0}_{\Bbb{Q}}(\ParL-\mathcal{R}_L)$ is surjective.
\item[(I2)] The map $\Pic^{+,\deg=0}_{\Bbb{Q}}(\ParL-\mathcal{R}_L)\to \Pic(\Par_{n,\mn,S})_{\Bbb{Q}}$
is injective with image exactly $\Pic'_{\Bbb{Q}}$.
\item[(I3)]  $\Pic^{+,\deg=0}_{\Bbb{Q}}(\ParL)$ is identified with  $\Pic^+_{\Bbb{Q}}(\Par_{r,\mathcal{O}(-d),S})\times \Pic^+_{\Bbb{Q}}(\Par_{n-r,\mathcal{O}(d-D),S})$,
\end{enumerate}
Putting these together one gets the desired surjection \eqref{inductiones}. Explicit formulas for the induction morphism will be given in Section \ref{explicitF}.

\subsection{Construction of the induction morphism}
If $1\not\in I^i$ for all $i$, we have a morphism  (see \eqref{morph1}) $\widehat{\Omega}^0(d,r,\mn,n,\vec{I})\to \ParL$ which fits into a diagram:
\begin{equation}\label{Dbasic1.5}
\source{rigid-local-systems-multiplicative-eigenvalue-problem}
\xymatrix{
 & \widehat{\Omega}^0(d,r,\mn,n,\vec{I})\ar[dl]^{\pi}\ar[dr]\\
\Par_{n,\mn,S}   &  & \ParL\ar@/_/[ul]_i\ar[ll]^{i'}
}
\end{equation}
Suppose an $s$-fold application of shifting (at various points of $S$, perhaps repeated) leads to $1\not\in I^i$ for all $i$. Note that there is a minimal choice of these shift operations, and we will use this choice (note that all shifts will turn out to produce the same induction map).

Let the shifted data be $(d',r,\mn',n,\vec{K})$. Set $D'=-\deg(\mn')$ and let $\op{Sh}$ denote the composites $\ParL(-d,r,n-r,d-D)\to \ParL(-d',r,n-r,d'-D')$, as well as $\Par_{n,\mn,S}\to \Par_{n,\mn',S}$(see the diagram \eqref{basicDnew}). To define the induction map, we take line bundles on  $\ParL(-d,r,n-r,d-D)$, obtain line bundles on $\ParL(-d',r,n-r,d'-D')$ by the inverse of shifting procedure $\Sh$. These line bundles then give line bundles on $\widehat{\Omega}^0(d',r,\mn',n,\vec{K})$, hence on $\widehat{\Omega}^0(d',r,\mn',n,\vec{K})-\mathcal{R}'$, where $\mathcal{R}'$ is the ramification locus of $\widehat{\Omega}^0(d',r,\mn',n,\vec{K})\to \Par_{n,\mn',S}$. By Corollary \ref{carole}, $\widehat{\Omega}^0(d',r,\mn',n,\vec{K})-\mathcal{R}'$ is an open substack of $\Par_{n,\mn',S}$ whose complement is of codimension $\geq 2$. Therefore we obtain a line bundle on
$\Par_{n,\mn',S}$  and we pull these line bundles to $\Par_{n,\mn,S}$ via $\Sh$.
\subsection{Proof of property (I1)}

\begin{lemma}
\source{rigid-local-systems-multiplicative-eigenvalue-problem}
\notready\label{prius2}
\source{rigid-local-systems-multiplicative-eigenvalue-problem}
The restriction mapping $\Pic(\ParL)\to \Pic(\ParL -\mathcal{R}_L)$ has a section and is hence surjective.
\end{lemma}
\begin{proof}
We may assume that $1\not\in I^i$ for all $i$ without loss of generality. Let $\mathcal{M}\in\Pic(\ParL -\mathcal{R}_L)$. To define the section we may just restrict
$\op{Ind}(\mathcal{M})$  via $\pi\circ i:\ParL \to \Par_{n,\mn,S}$.
\end{proof}
Property (I1) is the following:
\begin{proposition}
\source{rigid-local-systems-multiplicative-eigenvalue-problem}
\notready\label{M11}
\source{rigid-local-systems-multiplicative-eigenvalue-problem}
The restriction map $\Pic^{+,\deg=0}(\ParL)\to \Pic^{+,\deg=0}(\ParL -\mathcal{R}_L)$ is a surjection.
\end{proposition}
\begin{proof}
Let $\mathcal{M}$ be an effective line bundle on $\ParL -\mathcal{R}_L$ with section $s$. Then $\op{Ind}(\mathcal{M})$ is defined as a pullback of $\mathcal{M}$ and   therefore has a non-zero section $\tilde{s}$ defined as a pullback of $s$. The restriction of a $\tilde{s}$  to $\ParL$ coincides with $s$ on $\ParL -\mathcal{R}_L$ and is hence non-zero. This proves the proposition.
\end{proof}



\subsection{Proof of property (I2)}

\begin{defi}
\source{rigid-local-systems-multiplicative-eigenvalue-problem}
\notready  When $1\not\in I^i$ for all $i$, we consider $\Par'_{n,\mn,S}$ the parameter space of rank $n$ parabolic bundles with $\gamma:\det \mw\leto{\sim}\mn$, where for each $i$,  $\me^{p_i}$ now parametrizes {\em partial flags} where only $E^{p_i}_a$ with $a\in T(I^i)$ are defined (see Definition  \ref{ToI}).
\end{defi}
Using  the method of \cite{LS}, we obtain
\begin{lemma}
\source{rigid-local-systems-multiplicative-eigenvalue-problem}
\notready
 Assume $1\not\in I^i$ for all $i$. There is a natural mapping $\Par_{n,\mn,S}\to \Par'_{n,\mn,S}$.
The subgroup of $\Pic(\Par_{n,\mn,S})$ formed by $\Pic(\Par'_{n,\mn,S})$ coincides with the subgroup
$\Pic'$ defined in Proposition \ref{grill}.
\end{lemma}

\begin{proposition}
\source{rigid-local-systems-multiplicative-eigenvalue-problem}
\notready\label{grill}
\source{rigid-local-systems-multiplicative-eigenvalue-problem}
The image of
$$\op{Ind}:\Pic(\ParL -\mathcal{R}_L)\to \Pic(\Par_{n,\mathcal{N},S})$$
 is contained in  $\Pic'$.
\end{proposition}
\begin{proof}(of Proposition \ref{grill})
Using the shift operations, we may assume that $1\not\in I^i$ for all $i$.
Some parts of the flag structures on $\widehat{\Omega}^0(d,r,\mn,n,\vec{I})$ are ``unnecessary" for its definition. In fact there is a natural  representable $\widehat{\Omega'}^0(d,r,\mn,n,\vec{I})\to \Par'_{n,\mn,S}$ such that the base change to $\Par_{n,\mn,S}\to\Par'_{n,\mn,S}$ gives $\widehat{\Omega}^0(d,r,\mn,n,\vec{I})\to \Par_{n,\mn,S}$. Note also that by \cite[Lemma 7.10]{BHermit}, we have a natural map $\Par'_{n,\mn,S}\to \ParL$ which factorizes $\Par_{n,\mn,S}\to \ParL$.

Therefore assuming $1\not\in I^i$ for all $i$, the induction map \eqref{inducto} factors through  $\Pic(\Par'_{n,\mn,S})$ as desired.
%$$\op{Ind}:\Pic(\ParL-\mathcal{R}_L)\to \Pic(\Par'_{n,\mn,S})\subseteq \Pic(\Par_{n,\mn,S}).$$
\end{proof}


 \subsubsection{The mapping  \eqref{inducto} is an injection with image $\Pic'$}
\begin{proposition}
\source{rigid-local-systems-multiplicative-eigenvalue-problem}
\notready\label{outside}
\source{rigid-local-systems-multiplicative-eigenvalue-problem}
The induction morphism \eqref{inducto} is a bijection of $\Pic(\ParL -\mathcal{R}_L)$ with $\Pic'$.
\end{proposition}

Again to prove this proposition we may assume  $1\not\in I^i$ for all $i$. It follows immediately from the two lemmas below.

\begin{lemma}
\source{rigid-local-systems-multiplicative-eigenvalue-problem}
\notready\label{later}
\source{rigid-local-systems-multiplicative-eigenvalue-problem}
Assume $1\not\in I^i$ for all $i$.
The composite
$$\Pic(\Par'_{n,\mn,S})\to \Pic(\Par_{n,\mn,S})\to  \Pic({\Omega}^0(d,r,\mn,n,\vec{I})-\mathcal{R})$$ is an isomorphism.
\end{lemma}
\begin{proof}
We will construct an inverse of the mapping $\Pic(\Par'_{n,\mn,S})\to  \Pic({\Omega}^0(d,r,\mn,n,\vec{I})-\mathcal{R})$. Take a line bundle on ${\Omega}^0(d,r,\mn,n,\vec{I})-\mathcal{R}$ and extend it to $\Par_{n,\mn,S}$. There are many choices here, any two choices differ by a linear combination of the $D(a,j)$. There is a unique such extension which is in $\Pic'(\Par'_{n,\mn,S})$ by Lemma \ref{oneoff} (that the failure is by $1$  there is crucial). This gives the desired inverse. The two composites are identity by easy inspection.
\end{proof}

\begin{lemma}
\source{rigid-local-systems-multiplicative-eigenvalue-problem}
\notready\label{muchlater}
\source{rigid-local-systems-multiplicative-eigenvalue-problem}
Assume $1\not\in I^i$ for all $i$,
The pullback map
$\Pic(\Par_L-\mathcal{R}_L)\to  \Pic({\Omega}^0(d,r,\mn,n,\vec{I})-\mathcal{R})$is an isomorphism.
\end{lemma}
\begin{proof}
The section $i$ in \eqref{Dbasic} gives a section $\ParL-\mathcal{R}_L\to {\Omega}^0(d,r,\mn,n,\vec{I})-\mathcal{R}$ of the pullback map (see the Proof of \eqref{fulC}). Therefore we need to show only that a line bundle $\ma$ on ${\Omega}^0(d,r,\mn,n,\vec{I})-\mathcal{R}$ which restricts via $i$ to the trivial bundle
on $\ParL-\mathcal{R}_L$ is itself trivial. We wish to propagate the non-zero section of $\ma$ on  $\ParL-\mathcal{R}_L$ to a non-zero section on all of ${\Omega}^0(d,r,\mn,n,\vec{I})-\mathcal{R}$.
By  the Levification map (\cite[Section 3.8]{BKq}) we construct a map:
\begin{itemize}
\item[-] A map $\phi:\Bbb{A}^1\times {\Omega}^0(d,r,\mn,n,\vec{I})\to \Omega^0(d,r,\mn,n,\vec{I})$ such that $\phi_0$ equals $\Omega^0(d,r,\mn,n,\vec{I})\to \ParL$ followed by the section (where $\phi_t$ is the restriction of $\phi$ to $t\times \Omega^0(d,r,\mn,n,\vec{I})$).
\end{itemize}

We have a section of $\mathcal{A}$ on $\ParL$. Let $x\in \Omega^0(d,r,\mn,n,\vec{I})$. We obtain a line bundle on $\Bbb{A}^1$ by pulling back  $\mathcal{A}$ along $\Bbb{A}^1\to \Omega^0(d,r,\mn,n,\vec{I})$, $x\mapsto \phi(t,x)$. This line bundle is $\Bbb{G}_m$ equivariant, with trivial Mumford number at $0$. It is therefore canonically trivial (since we have a given section at $t=0$). We therefore obtain a canonical element in $\mathcal{T}_x$ for any $x$. This shows $\mathcal{T}$ is trivial (see \cite[Lemma 8.6]{BHermit})).

If  $\mathcal{T}\in \Pic'(\mf)^{\deg=0}$, then it is easy to see that $\mathcal{M}$ is in $\Pic^{\deg=0}(\ParL -\mathcal{R}_L)$, since the degree on both sides is the weight of the action of the center of the Levi subgroup in both cases.
\end{proof}
\subsubsection{Proof of property (I2)}
Property (I2) follows from Proposition \ref{outside} and  the following:
\begin{proposition}
\source{rigid-local-systems-multiplicative-eigenvalue-problem}
\notready\label{uselater}
\source{rigid-local-systems-multiplicative-eigenvalue-problem}
Suppose $\mathcal{M}\in \Pic^{\deg=0}(\ParL-\mathcal{R}_L)$. Then
\begin{enumerate}
\item[(a)] $\op{Ind}(\mathcal{M})\in \Pic^{\deg=0}(\Par_{n,\mn,S})$.
\item[(b)] $H^0(\ParL-\mathcal{R}_L,\mathcal{M})\to H^0(\Par_{n,\mn,S},\op{Ind}(\mathcal{M}))$ is an isomorphism
\item[(c)] $H^0(\ParL,\op{Ind}(\mathcal{M})|_{\ParL})\to H^0(\ParL -\mathcal{R}_L,\mathcal{M})$ is also an isomorphism.
\end{enumerate}
\end{proposition}
\begin{proof}
For part (a), let $\op{Ind}(\mathcal{M})=\mathcal{B}(\vec{\lambda},\ell)$. Now $\op{Ind}(\mathcal{M})$ restricted to $\ParL-\mathcal{R}_L$ lies in  $\Pic^{\deg=0}(\ParL-\mathcal{R}_L)$, i.e., has index zero.
The index is easily computed to be  up to a multiple equal to the difference between the two sides of \eqref{wallneww} (Note that multiplication by $t$ on a vector bundle acts on the determinant of cohomology by $t^{\chi}$, with $\chi$ the Euler characteristic of the vector bundle).  Part (a)  follows.

Part (b) is a generalization of Proposition  \ref{fulC} (for $\mathcal{M}=\mathcal{O}$), and the proof is similar. The map in (c) is obviously injective. Since we have  maps (reduce first to the case $1\not\in I^i$ for all $i$)
$$ H^0(\ParL -\mathcal{R}_L,\mathcal{M})\to H^0(\Par_{n,\mn,S},\op{Ind}(\mathcal{M}))\to H^0(\ParL,\op{Ind}(\mathcal{M})|_{\ParL})$$
the surjectivity follows.
\end{proof}
\subsection{Proof of property (I3)}
We have a natural map
\begin{equation}\label{temper2}
\source{rigid-local-systems-multiplicative-eigenvalue-problem}
{\Par}_{r,\mathcal{O}(-d),S}\times {\Par}_{n-r,\mathcal{O}(d-D),S} \to  \ParL.
\end{equation}

\begin{lemma}
\source{rigid-local-systems-multiplicative-eigenvalue-problem}
\notready\label{BachMass}
\source{rigid-local-systems-multiplicative-eigenvalue-problem}
\begin{enumerate}
\item  $\Pic^{+,\deg=0}_{\Bbb{Q}}(\ParL)\leto{\sim}\Pic^+_{\Bbb{Q}}(\Par_{r,\mathcal{O}(-d),S})\times \Pic^+_{\Bbb{Q}}(\Par_{n-r,\mathcal{O}(d),S})$.
\item $\Pic^{+}(\ParL)\leto{\sim}\Pic^+(\Par_{r,\mathcal{O}(-d),S})\times \Pic^+(\Par_{n-r,\mathcal{O}(d),S})$.
\item Under the correspondence in (2), a line bundle $\ml$ on $\ParL$ corresponds to tensor product $\ml_1\boxtimes\ml_2$ with one of $\ml_1$ a F-line bundle and the other trivial, if and only if,
$\ml=\mathcal{O}(E)$ with $E$ non-empty, reduced and irreducible.

\end{enumerate}
\end{lemma}
This result follows from Lemma \ref{mass2} below applied to \eqref{temper2}.
It is necessary to take $\Bbb{Q}$-coefficients above for the reverse map in (1).






\subsection{Induction and extremal rays coming from F-line bundles}
We want to determine the set of F-line bundles in $\mf^{(2)}$. Note that using shift operations (Proposition \ref{easily}), we can put the set of F-line bundles on $\Par_{r,\mathcal{N},S}$ in bijection with the set of F-line bundles
on   $\Par_{r,\mathcal{O},S}$ (which are in bijection with F-vertices of $P_r(s)$ by Lemma \ref{corresp}).
\begin{proposition}
\source{rigid-local-systems-multiplicative-eigenvalue-problem}
\notready\label{condensed}
\source{rigid-local-systems-multiplicative-eigenvalue-problem}
There is a bijection between the following two sets:
\begin{enumerate}
\item The set of  $\ml\in \Pic'\subseteq\Pic(\Par_{n,\mathcal{O},S})$ (see Definition  \ref{Br2}) such that $\ml$ is a F-line bundle on  $\Par_{n,\mathcal{O},S}$.
\item The set of $\mathcal{M}\in  \Pic^{+}(\ParL)$ such that the following two conditions hold
\begin{enumerate}
\item[(a)]$\mathcal{M}$ corresponds to a F-line bundle on one of the factors of (and $0$ on the other)
$$\Pic^+(\Par_{r,\mathcal{O}(-d),S})\times \Pic^+(\Par_{n-r,\mathcal{O}(d),S})\subseteq \Pic^+_{\Bbb{Q}}(\Par_{r,\mathcal{O}(-d),S})\times \Pic^+_{\Bbb{Q}}(\Par_{n-r,\mathcal{O}(d),S}) $$ under the  bijection in Lemma \ref{BachMass}.
\item[(b)] $Ind(\mathcal{M})\neq 0$ and $H^0(\ParL,\op{Ind}(\mathcal{M})|_{\ParL})$ (which breaks up as a tensor product)  is one-dimensional (see Remark \ref{putR} for equivalent conditions).
\end{enumerate}
\end{enumerate}
The bijection takes $\mathcal{M}$ in (2) to its induction $\ml=\op{Ind}(\mathcal{M})$.
\end{proposition}

\begin{proof}
We start by showing that if $\mathcal{M}=\mathcal{O}(D)$ is as in (2), then $\ml=\op{Ind}(\mathcal{M})$ is an F-line bundle.
By Proposition \ref{uselater}, and (b), $H^0(\Par_{n,\mathcal{O},S},\ml)$ is one dimensional. Let $\ml=
\mathcal{O}(E)$ and we want to show $E$ is irreducible and reduced. Suppose $E=E_1+E_2$, let $\ml_1=\mathcal{O}(E_1)$ and $\ml_2=\mathcal{O}(E_2)$. There exist $\mathcal{M}_1$ and $\mathcal{M}_2$ effective line bundles on $\Par_L$ which induce to $\ml_1$ and $\ml_2$. Let $\mathcal{M}_i=\mathcal{O}(D_i)$. By Proposition \ref{uselater}, $D_1+D_2$ and $D$ agree on $\Par_L-\mathcal{R}_L$. Since $D$ is irreducible, either $D_1$ or $D_2$ is supported on $\mathcal{R}_L$, so either $\ml_1$ or $\ml_2$ (being the inductions of $\mathcal{M}_1$ and $\mathcal{M}_2$) is trivial, as desired. If both $D_1$ and $D_2$ are both supported on  $\mathcal{R}_L$, then $D$ is also supported on $\mathcal{R}_L$, and hence $\op{Ind}(\mathcal{M})$ is trivial.

To go the other direction, start with a $\ml=\mathcal{O}(E)$ as in (1). By Proposition \ref{uselater}, $H^0(\Par_L-\mathcal{R}_L,\mathcal{L}|_{\ParL})$ is  one dimensional. Let ${D}\subset \ParL$ be the closure of the zeroes of this section and set $\mathcal{M}=\mathcal{O}(D)$.  We claim that $\mathcal{M}$ satisfies the conditions in (2). By construction $\op{Ind}(\mathcal{M})=\ml$, and hence $\mathcal{M}$ is non trivial. The condition in (b) holds by Proposition \ref{uselater}.

We now show that $\mathcal{M}$ is a F-line bundle. Since
$\mathcal{L}|_{\ParL}$ and $\mathcal{M}$ agree on $\Par_L-\mathcal{R}_L$, $H^0(\Par_L,\mathcal{M})$ is one dimensional.


If $D=D_1+D_2$, then we get a corresponding decomposition of $\ml=\ml_1\tensor\ml_2$ as a tensor product of effective bundles. We see that one of the factors $\ml_1$ or $\ml_2$ is trivial, and hence $D_1$ or $D_2$ is trivial on $\Par_L-\mathcal{R}_L$. Since $D$ is the closure of its intersection with $\Par_L-\mathcal{R}_L$, this makes one of $D_1$ and  $D_2$ trivial, and hence $\mathcal{M}$ is an F-line bundle.

The composite of going from (2) to (1) and back to (2) as described above is the identity. This finishes the proof of the proposition.


\end{proof}
\begin{remark}
\source{rigid-local-systems-multiplicative-eigenvalue-problem}
\notready\label{putR}
\source{rigid-local-systems-multiplicative-eigenvalue-problem}
The condition in (b) is equivalent to $Ind(\mathcal{M})\neq 0$ and $h^0(\ParL,\mathcal{M}(\mathcal{R}_L))=1$. This is because any section of $H^0(\ParL-\mathcal{R}_L,\mathcal{M})$ lifts to a section of
$H^0(\ParL,\mathcal{M}(m\mathcal{R}_L))$ for some $m>0$, and we know that $\mathcal{M}$ is effective and $h^0(\ParL,\mathcal{M}^m(m\mathcal{R}_L))=1$ by quantum Fulton.

\end{remark}
